\chapter{Convex Functions}
\label{chap:ch01}

% Chapter goal: build intuition for convexity (a bowl with no false
% bottoms) and show why it is the dividing line between optimization
% problems that are easy and those that are hard.

\section{Why Convexity Matters for AI}
\label{sec:ch01-why}

\begin{aibox}
  Linear regression, logistic regression, and SVMs have convex losses, so
  gradient descent finds the global minimum. Deep networks do not, and much
  of modern training practice exists to work around that.
\end{aibox}

\section{Functions: A Quick Refresher}
\label{sec:ch01-functions}
\subsection{Domain, Range, and Graphs}
\subsection{Functions You Will See Everywhere in AI}
% linear/affine, quadratic, exp/log, sigmoid, ReLU, softplus, log-sum-exp

\section{Convex Sets}
\label{sec:ch01-convex-sets}
\subsection{Line Segments and Convex Combinations}
\subsection{Examples: Lines, Half-Spaces, Balls, the Probability Simplex}

\section{Definition of a Convex Function}
\label{sec:ch01-definition}

\begin{definition}[Convex function]
  \label{def:convex}
  A function $f\colon \dom f \to \R$, where $\dom f$ is a convex set, is
  \emph{convex} if for all $x, y \in \dom f$ and all $\theta \in [0,1]$,
  \begin{equation}
    f\bigl(\theta x + (1-\theta) y\bigr) \le \theta f(x) + (1-\theta) f(y).
    \label{eq:convex}
  \end{equation}
\end{definition}

\subsection{Geometric Picture: The Chord Lies Above the Graph}
\subsection{Strict and Strong Convexity}
\subsection{Concave Functions}
\subsection{The Epigraph View}

\section{Examples and Non-Examples}
\label{sec:ch01-examples}
\subsection{Convex: Affine, Quadratic, Exponential, Norms, Log-Sum-Exp}
\subsection{Not Convex: Sine, Cubic, Neural Network Losses}

\section{Recognizing Convexity with Derivatives}
\label{sec:ch01-conditions}
% Light preview; derivatives are developed fully in Chapter 2.
\subsection{First-Order Condition: Tangent Lines Lie Below the Graph}
\subsection{Second-Order Condition: Non-Negative Curvature}

\section{Operations That Preserve Convexity}
\label{sec:ch01-operations}
\subsection{Non-Negative Weighted Sums}
\subsection{Composition with an Affine Map}
\subsection{Pointwise Maximum (and Why ReLU and Hinge Loss Are Convex)}

\section{Jensen's Inequality}
\label{sec:ch01-jensen}

\section{Why Convex Problems Are Easy: Local Minima Are Global}
\label{sec:ch01-local-global}

\section{Convexity in Machine Learning}
\label{sec:ch01-ml}
\subsection{Mean Squared Error and Linear Regression}
\subsection{Cross-Entropy and Logistic Regression}
\subsection{Hinge Loss and Support Vector Machines}
\subsection{Regularizers: $L_1$ and $L_2$ Penalties}
\subsection{Beyond Convexity: The Loss Landscape of Deep Networks}

\section{In Code: Checking and Visualizing Convexity}
\label{sec:ch01-code}
% Plot functions with chords, numerically test inequality (1.1), and run
% gradient descent on a convex vs. a non-convex function.

\section{Summary}

\section{Exercises}

\section{Further Reading}
Boyd and Vandenberghe~\cite{boyd2004convex} is the standard reference on
convex functions and convex optimization.
